
\documentclass{article}
\usepackage{colortbl}
\usepackage{makecell}
\usepackage{multirow}
\usepackage{supertabular}

\begin{document}

\newcounter{utterance}

\centering \large Interaction Transcript for game `cladder', experiment `full\_v1.5\_default', episode 2646 with selfrefine{-}Qwen3.5{-}27B.
\vspace{24pt}

{ \footnotesize  \setcounter{utterance}{1}
\setlength{\tabcolsep}{0pt}
\begin{supertabular}{c@{$\;$}|p{.15\linewidth}@{}p{.15\linewidth}p{.15\linewidth}p{.15\linewidth}p{.15\linewidth}p{.15\linewidth}}
    \# & \multicolumn{2}{c}{Player} && \multicolumn{2}{c}{Game Master} \\
    \hline

    \theutterance \stepcounter{utterance}  
    & & & \multicolumn{4}{p{0.6\linewidth}}{
        \cellcolor[rgb]{0.9,0.9,0.9}{
            \makecell[{{p{\linewidth}}}]{
                \texttt{\tiny{[P1$\langle$GM]}}
                \texttt{You are an expert in causal inference. The following question is not a typical commonsense query, but rather a meticulously designed question created by a professor specializing in causal inference, intended to assess the students' mastery of the course content.} \\
\texttt{Imagine a self{-}contained, hypothetical world with only the following conditions, and without any unmentioned factors or causal relationships: Gender has a direct effect on smoking and lung cancer. Smoking has a direct effect on tar deposit. Tar deposit has a direct effect on lung cancer. Gender is unobserved. For nonsmokers and with no tar deposit, the probability of lung cancer is 79\%. For nonsmokers and with high tar deposit, the probability of lung cancer is 49\%. For smokers and with no tar deposit, the probability of lung cancer is 71\%. For smokers and with high tar deposit, the probability of lung cancer is 40\%. For nonsmokers, the probability of high tar deposit is 79\%. For smokers, the probability of high tar deposit is 40\%. For smokers, would it be more likely to see lung cancer if the person had been a nonsmoker?} \\
\texttt{Start your answer with 'Yes' or 'No', followed by additional reasoning or evidence to support your explanation.} \\
            }
        }
    }
     \\ \\

    \theutterance \stepcounter{utterance}  
    & \multicolumn{4}{p{0.6\linewidth}}{
        \cellcolor[rgb]{0.561,0.737,0.561}{
            \makecell[{{p{\linewidth}}}]{
                \texttt{\tiny{[P1$\rangle$GM]}}
                \texttt{No} \\
\\ 
\texttt{To determine whether a smoker would be more likely to see lung cancer if they had been a nonsmoker, we must calculate the counterfactual probability of lung cancer under the intervention \$do(S=0)\$ (making the person a nonsmoker) and compare it to the observed probability of lung cancer for smokers (\$S=1\$).} \\
\\ 
\texttt{In this hypothetical world, the causal structure is defined as:} \\
\texttt{1.  \$G \textbackslash{}to S\$ (Gender affects Smoking)} \\
\texttt{2.  \$G \textbackslash{}to L\$ (Gender affects Lung Cancer)} \\
\texttt{3.  \$S \textbackslash{}to T\$ (Smoking affects Tar)} \\
\texttt{4.  \$T \textbackslash{}to L\$ (Tar affects Lung Cancer)} \\
\\ 
\texttt{Since Gender (\$G\$) is unobserved, we cannot simply condition on it to block the backdoor path \$S \textbackslash{}leftarrow G \textbackslash{}to L\$. However, we are given the full conditional probability distributions for the observed variables (\$S\$ and \$T\$) and the outcome (\$L\$). We can calculate the expected probability of lung cancer under the intervention \$do(S=0)\$ and compare it to the observed probability for smokers.} \\
\\ 
\texttt{**Step 1: Calculate the observed probability of Lung Cancer for Smokers (\$P(L=1 | S=1)\$)**} \\
\texttt{We use the law of total probability, conditioning on Tar (\$T\$):} \\
\texttt{\$\$P(L=1 | S=1) = P(L=1 | S=1, T=0)P(T=0 | S=1) + P(L=1 | S=1, T=1)P(T=1 | S=1)\$\$} \\
\\ 
\texttt{Given data:} \\
\texttt{*   \$P(L=1 | S=1, T=0) = 0.71\$} \\
\texttt{*   \$P(L=1 | S=1, T=1) = 0.40\$} \\
\texttt{*   \$P(T=1 | S=1) = 0.40\$ (Probability of high tar for smokers)} \\
\texttt{*   \$P(T=0 | S=1) = 1 {-} 0.40 = 0.60\$} \\
\\ 
\texttt{Calculation:} \\
\texttt{\$\$P(L=1 | S=1) = (0.71 \textbackslash{}times 0.60) + (0.40 \textbackslash{}times 0.40)\$\$} \\
\texttt{\$\$P(L=1 | S=1) = 0.426 + 0.160 = 0.586\$\$} \\
\\ 
\texttt{**Step 2: Calculate the counterfactual probability of Lung Cancer if the smoker were a nonsmoker (\$P(L=1 | do(S=0))\$)**} \\
\texttt{Under the intervention \$do(S=0)\$, the variable \$S\$ is set to 0 (nonsmoker). This breaks the causal link from \$G\$ to \$S} \\
            }
        }
    }
    & & \\ \\

    \theutterance \stepcounter{utterance}  
    & & & \multicolumn{2}{p{0.3\linewidth}}{
        \cellcolor[rgb]{0.9,0.9,0.9}{
            \makecell[{{p{\linewidth}}}]{
                \texttt{\tiny{[GM$|$GM]}}
                \texttt{No} \\
            }
        }
    }
    & & \\ \\

    \theutterance \stepcounter{utterance}  
    & & & \multicolumn{2}{p{0.3\linewidth}}{
        \cellcolor[rgb]{0.9,0.9,0.9}{
            \makecell[{{p{\linewidth}}}]{
                \texttt{\tiny{[GM$|$GM]}}
                \texttt{no} \\
            }
        }
    }
    & & \\ \\

    \theutterance \stepcounter{utterance}  
    & & & \multicolumn{2}{p{0.3\linewidth}}{
        \cellcolor[rgb]{0.9,0.9,0.9}{
            \makecell[{{p{\linewidth}}}]{
                \texttt{\tiny{[GM$|$GM]}}
                \texttt{game\_result = WIN} \\
            }
        }
    }
    & & \\ \\

\end{supertabular}
}

\end{document}
